\documentclass[11pt]{article}
\usepackage[margin=2.4cm]{geometry}
\usepackage{fontspec}
\setmainfont{DejaVu Sans}
\usepackage{juliamono}
\usepackage{fvextra}
\usepackage{xcolor}
\usepackage{xurl}
\usepackage{hyperref}
\usepackage{amssymb}
\newcommand{\checkbox}{$\square$}
\newcommand{\answerchoices}{\checkbox\ YES \quad \checkbox\ NO \quad \checkbox\ CANNOT DETERMINE}
\newcommand{\decisionexplanation}{%
  \par\smallskip\textbf{If NO, give the reason for the mismatch (required). If
  CANNOT DETERMINE, state what context or evidence is missing:}\par
  \rule{\linewidth}{0.4pt}\par\vspace{7pt}\rule{\linewidth}{0.4pt}\par}
\setlength{\parindent}{0pt}\setlength{\parskip}{5pt}
\setlength{\emergencystretch}{2em}
\DefineVerbatimEnvironment{Vsrc}{Verbatim}{fontsize=\small,frame=single,
  breaklines=true,breakanywhere=true}
\DefineVerbatimEnvironment{Vlean}{Verbatim}{fontsize=\scriptsize,frame=single,
  numbers=left,numbersep=5pt,breaklines=true,breakanywhere=true}
\hypersetup{colorlinks=true,urlcolor=blue}

\begin{document}
\begin{center}{\LARGE Proof Verification Document}\\[2pt]{\large reddit\_learnmath\_t6wb5n}\\[2pt]{\small Reddit proof-check posts (newly confirmed)}\end{center}
\begin{center}\fcolorbox{blue!60!black}{blue!4}{\parbox{0.92\linewidth}{\textbf{Make three independent decisions.} Follow the evidence chain: original source $\longrightarrow$ challenged claim $\longrightarrow$ proposed counterexample $\longrightarrow$ Lean theorem. A record is valid only when all three links are correct. Complete the human sign-off before reading the automated review appendix; those earlier results are background, not an answer key.}}\end{center}
\textbf{Use only this document, except for the public source link.} Everything needed for Questions 2 and 3 is included here. For Question 1, open the public source link when one is shown. If there is no public link, the complete stored problem and solution are included here. If the information is still not enough, choose CANNOT DETERMINE.
For each question, choose \textbf{YES}, \textbf{NO}, or \textbf{CANNOT DETERMINE}. Choose CANNOT DETERMINE if the document does not provide enough information for you to decide; do not mark YES if you are uncertain. If you find a specific mismatch, choose NO and explain it in the space below that question. Do not rewrite the claim, add assumptions, or substitute a different counterexample.
\textbf{How copied text is shown.} The problem, quoted passage, and proposed counterexample come from the stored record. The Lean theorem statement and Lean files are copied exactly from the certificate. Some long lines continue on the next printed line; this is only page formatting and does not change the copied text. Line numbers are added only to help you find passages. If the stored record contains a character that cannot be printed, the document shows \texttt{[U+XXXX CONTROL]} in its place.
\begin{description}
\item[Challenged claim.] The exact statement from the source that this record says is false.
\item[Proposed counterexample.] A concrete case that should satisfy the claim's assumptions but make its conclusion false.
\item[Lean certificate.] A Lean file proving a formal theorem. Successful compilation proves the theorem as written, not that it is the right translation of the source.
\end{description}
\section*{Question 1 --- Does this proof verification document accurately represent the source?}
\textbf{What to decide.} Compare the challenged passage below with the original source---either the public link or the complete source text in this document. Answer YES only if they make the same claim and include the same assumptions or conditions. If this document changes the meaning, leaves out a condition, or points to a different passage, answer NO.
\subsection*{Original source}
\textbf{Source record:}
\begin{Vsrc}
qedbench_certificate_availability_triage/data/records_reddit_exhaustive.jsonl
\end{Vsrc}
\textbf{Public source:}\par
\url{https://www.reddit.com/r/learnmath/comments/t6wb5n/rational_epsilondelta_limit_proof_questions/}\par
\subsection*{Problem or question from the source (verbatim)}
\begin{Vsrc}
Rational Epsilon-Delta Limit Proof Questions
\end{Vsrc}
\subsection*{Challenged passage recorded from the source (stored text)}
\begin{Vsrc}
|(1/x\^2)-(1/4)| = |(4-x\^2)/(4x\^2)| = |4-x\^2|/|4x\^2| = |x\^2-4|/(4x\^2) = |(x-2)(x+2)|/(4x\^2) = |x-2||x+2|/(x\^2).
\end{Vsrc}
\subsection*{Surrounding source lines (stored text with added line labels)}
The labels refer to lines in the solution text stored with the source record; the source-line content is not transliterated.
\begin{Vsrc}
0001 | Good afternoon. I have more questions about the details of epsilon-delta proofs. Below is a simple, rational limit proof example with questions at the end. The scratch work and proof are a bit pedantic but I don't follow proofs very well which omit a lot of details, including scratch work and thinking process.
0002 | 
0003 | Claim: lim (x->2) (1/x\^2) = (1/4).
0004 | 
0005 | WTS: For every epsilon > 0, there exists a delta greater than zero such that, when x is any real number, then 0<|x-2|<delta => |(1/x\^2)-(1/4)|<epsilon.
0006 | 
0007 | &#x200B;
0008 | 
0009 | Scratch Work:
0010 | 
0011 | \* Manipulate implication 0<|x-2|<delta => |(1/x\^2)-(1/4)|<epsilon to find delta.
0012 | 
0013 | \* So |(1/x\^2)-(1/4)| = |(4-x\^2)/(4x\^2)| = |4-x\^2|/|4x\^2| = |x\^2-4|/(4x\^2) = |(x-2)(x+2)|/(4x\^2) = |x-2||x+2|/(x\^2).
0014 | 
0015 | \* By assumption, an upper bound on |x-2| is delta.
0016 | 
0017 | \* Need to find an upper bound on |x+2|/(x\^2) by making |x+2|/(x\^2) < (C/4\*D) for some numbers C and D. Then any delta <=(4\*D/C)\*epsilon will bound |x+2|/(x\^2) above.
0018 | 
0019 | \* Choose delta <= 1.
0020 | 
0021 | \* Then, to find C: |x-2|<delta => -1<x-2<1 => 1<x<3 => 3<x+2<5 => -5<3<x+2<5 => |x+2|<5.
0022 | 
0023 | \* Similarly, to find D: |x-2|<delta => -1<x-2<1 => 1<x<3 => 1>(1/x)>(1/3) => 1>(1/x\^2)>(1/9) => (1/9)<(1/x\^2)<1 => -1<(1/9)<(1/x\^2)<1 => 1/|x\^2|<1 => 1/(x\^2)<1.
0024 | 
0025 | \* Thus (C/4\*D) = (5/4\*1) = (5/4).
\end{Vsrc}
\textbf{A YES means:} the proof verification document faithfully represents both the source context and the challenged claim, without dropping a condition or changing its meaning.
\section*{Question 2 --- Does the proposed counterexample disprove that claim?}
\textbf{Claim to test (repeated for convenience).}
\begin{Vsrc}
|(1/x\^2)-(1/4)| = |(4-x\^2)/(4x\^2)| = |4-x\^2|/|4x\^2| = |x\^2-4|/(4x\^2) = |(x-2)(x+2)|/(4x\^2) = |x-2||x+2|/(x\^2).
\end{Vsrc}
\subsection*{How the record says the argument uses this claim}
This is the record's description, not an established fact; check it against the source.
\begin{Vsrc}
The final equality |(x-2)(x+2)|/(4x^2) = |x-2||x+2|/(x^2) in line 0013 is valid.
\end{Vsrc}
\subsection*{Proposed counterexample (exact values or construction)}
\begin{Vsrc}
Take x=1, which satisfies the required domain condition x≠0. Then |(x−2)(x+2)|/(4x²)=|(-1)(3)|/4=3/4, while |x−2||x+2|/x²=|-1||3|=3. Since 3/4≠3, the claimed equality is false.
\end{Vsrc}
\subsection*{Why the record says the counterexample works}
Check this explanation independently; it is evidence to assess, not an answer key.
\begin{Vsrc}
The numerator identity |(x−2)(x+2)|=|x−2||x+2| is valid, but it does not justify dropping the factor 4 from the denominator.
\end{Vsrc}
\subsection*{How to decide}
\begin{enumerate}
\item List every assumption that applies to the challenged claim.
\item Check that the proposed example satisfies every one of those assumptions.
\item Recompute the claimed outcome and the actual outcome using exact reasoning.
\item Confirm that the conflict is with the challenged claim itself, not an unrelated sentence.
\end{enumerate}
\textbf{A YES means:} the example is fully specified, satisfies all relevant assumptions, and makes the challenged claim false. If a necessary value, condition, or calculation is missing, choose CANNOT DETERMINE.
\section*{Question 3 --- Does the final Lean theorem statement verify this same counterexample?}
\textbf{What to decide.} A Lean file can compile while proving the wrong theorem. Compare these three links independently:
\begin{enumerate}
\item The Lean definitions represent the same mathematical objects, assumptions, and operations as the challenged claim.
\item The Lean construction uses the same concrete values or example shown in Question 2.
\item The final theorem statement asserts that the challenged claim fails on that example, not merely a weaker or nearby statement.
\end{enumerate}
\subsection*{Selected final Lean theorem statement (verbatim UTF-8 excerpt)}
This is an exact contiguous slice of the certificate file, from the theorem keyword through the original \texttt{:=} token where its proof begins. It is shown first so the formal target is not buried in the complete file.
\begin{Vlean}
theorem reddit_learnmath_t6wb5n_refutation :
    nonzero witnessX /\
    nonzero (sourceLeftDenominator witnessX) /\
    nonzero (sourceRightDenominator witnessX) /\
    equivalent (sourceNumerator witnessX) (claimedNumerator witnessX) /\
    equivalent witnessLeft threeFourths /\
    equivalent witnessRight three /\
    Not (equivalent witnessLeft witnessRight) /\
    Not PublishedFinalEqualityAtWitness :=
\end{Vlean}
\subsection*{Project-local imported Lean files (complete, verbatim UTF-8)}
This section follows the certificate's import commands recursively and prints every repository-owned Lean source file on which the final theorem statement depends. Each file is copied in full, including its definitions and controls. External Mathlib, Std, and Lean Init source is not duplicated.
No project-local imported Lean files were found. External Mathlib, Std, and Lean Init source is not duplicated in this proof verification document.
\subsection*{Complete counterexample Lean file (verbatim UTF-8)}
The complete Lean file is included below so you can check every definition and supporting result used by the final theorem statement. Use it together with the project-local imported Lean files above.
\textbf{Lean certificate file:}
\begin{Vsrc}
partB_feasibility/reddit_certificates/reddit_learnmath_t6wb5n.lean
\end{Vsrc}
\begin{Vlean}
/-!
Record: reddit_learnmath_t6wb5n
Public source: https://www.reddit.com/r/learnmath/comments/t6wb5n/rational_epsilondelta_limit_proof_questions/

Exact source claim under review (the final equality in line 0013):
`|(x-2)(x+2)|/(4*x^2) = |x-2||x+2|/(x^2)`.
The source displayed it as the last step of
`|(1/x^2)-(1/4)| = |(4-x^2)/(4x^2)| = |4-x^2|/|4x^2|
 = |x^2-4|/(4x^2) = |(x-2)(x+2)|/(4x^2)
 = |x-2||x+2|/(x^2)`.

The left definition below deliberately retains the explicit factor `4`.
The numerator identity `|(x-2)(x+2)| = |x-2||x+2|` is valid and is
certified separately; it cannot erase that denominator factor.

Concrete independently verified witness: `x = 1`, with `x != 0`. Both
divisors are nonzero, the numerator identity holds, but the two source sides
are exactly `3/4` and `3`, so the final displayed equality is false.

A context-strengthened witness is also certified at `x = 3/2`,
`epsilon = 1`, and `delta = 4/5`: epsilon and delta are positive,
`delta = min(1,(4/5)*epsilon)`, `0 < |x-2| < delta`, `x != 0`, both
divisions are defined, the numerator identity holds, and the source sides are
`7/36` and `7/9`. Thus the step also fails inside the surrounding proof
neighborhood. Those neighborhood assumptions are not load-bearing for the
standalone algebraic equality, and repairability of the limit proof does not
make this local equality valid.

Negative control: `x = 2` is nonzero and both divisions are defined, but the
zero numerator makes both sides zero and masks the missing factor. This
detects a uniformly-false or vacuous encoding.

Route review and modelling judgment: `ExactRat` is an infinite, exact signed
rational model with an `Int` numerator, a positive `Nat` denominator, and
rational equivalence by cross multiplication rather than structure equality.
It implements integer constants, exact fractions, addition/subtraction,
multiplication, squaring, ordinary absolute value through `Int.natAbs`, and
division through a reciprocal whose interface requires a nonzero proof.

This model represents only the listed exact rational real witnesses. It is
not a finite-domain proxy and does not replace the source claim with a toy
finite analogue or a universal real-analysis, limit, topology, continuity, or
approximation claim. Exact rational field operations have the same values at
these rational real witnesses, so no imported real-analysis semantics are
needed. All results below are reducible kernel computations: there is no
floating point, bounded enumeration, imported tactic, or native evaluation.
-/

namespace RedditLearnmathT6wb5n

structure ExactRat where
  num : Int
  den : Nat
  den_pos : 0 < den

/- Lean's zero-import prelude does not provide the `DecidableEq Int` instance
used implicitly by the first certificate. These total Boolean comparisons
inspect the two core `Int` constructors directly. -/
def intEqual : Int -> Int -> Bool
  | Int.ofNat m, Int.ofNat n => m == n
  | Int.negSucc m, Int.negSucc n => m == n
  | _, _ => false

def intLess : Int -> Int -> Bool
  | Int.ofNat m, Int.ofNat n => decide (m < n)
  | Int.ofNat _, Int.negSucc _ => false
  | Int.negSucc _, Int.ofNat _ => true
  | Int.negSucc m, Int.negSucc n => decide (n < m)

/- Rational equivalence remains literal cross multiplication. These
Prop-valued computational predicates are abbreviations so zero-import kernel
`decide` can see their Boolean equalities while synthesizing `Decidable` in
Lean 4.30. -/
abbrev equivalent (a b : ExactRat) : Prop :=
  intEqual (a.num * Int.ofNat b.den) (b.num * Int.ofNat a.den) = true

abbrev nonzero (a : ExactRat) : Prop :=
  intEqual a.num 0 = false

def ofInt (n : Int) : ExactRat :=
  { num := n, den := 1, den_pos := by decide }

def fraction (n : Int) (d : Nat) (hd : 0 < d) : ExactRat :=
  { num := n, den := d, den_pos := hd }

def add (a b : ExactRat) : ExactRat :=
  { num := a.num * Int.ofNat b.den + b.num * Int.ofNat a.den
    den := a.den * b.den
    den_pos := Nat.mul_pos a.den_pos b.den_pos }

def sub (a b : ExactRat) : ExactRat :=
  { num := a.num * Int.ofNat b.den - b.num * Int.ofNat a.den
    den := a.den * b.den
    den_pos := Nat.mul_pos a.den_pos b.den_pos }

def mul (a b : ExactRat) : ExactRat :=
  { num := a.num * b.num
    den := a.den * b.den
    den_pos := Nat.mul_pos a.den_pos b.den_pos }

def square (a : ExactRat) : ExactRat :=
  mul a a

def absRat (a : ExactRat) : ExactRat :=
  { num := Int.ofNat a.num.natAbs
    den := a.den
    den_pos := a.den_pos }

/- The nonzero proof is a required input. Its impossible zero branch returns
a harmless total value; that branch cannot be called from a well-typed proof
of `nonzero a`. -/
def reciprocal (a : ExactRat) (_ha : nonzero a) : ExactRat :=
  match a.num with
  | Int.ofNat 0 =>
      { num := 0, den := 1, den_pos := by decide }
  | Int.ofNat (Nat.succ n) =>
      { num := Int.ofNat a.den
        den := Nat.succ n
        den_pos := Nat.zero_lt_succ n }
  | Int.negSucc n =>
      { num := -(Int.ofNat a.den)
        den := Nat.succ n
        den_pos := Nat.zero_lt_succ n }

def div (a b : ExactRat) (hb : nonzero b) : ExactRat :=
  mul a (reciprocal b hb)

abbrev less (a b : ExactRat) : Prop :=
  intLess (a.num * Int.ofNat b.den) (b.num * Int.ofNat a.den) = true

def lessOrEqual (a b : ExactRat) : Bool :=
  intLess (a.num * Int.ofNat b.den) (b.num * Int.ofNat a.den) ||
    intEqual (a.num * Int.ofNat b.den) (b.num * Int.ofNat a.den)

def minRat (a b : ExactRat) : ExactRat :=
  if lessOrEqual a b then a else b

/- These definitions preserve the source syntax. In particular,
`sourceLeftDenominator` contains the material factor `4`. -/
def sourceNumerator (x : ExactRat) : ExactRat :=
  absRat (mul (sub x (ofInt 2)) (add x (ofInt 2)))

def claimedNumerator (x : ExactRat) : ExactRat :=
  mul (absRat (sub x (ofInt 2))) (absRat (add x (ofInt 2)))

def sourceLeftDenominator (x : ExactRat) : ExactRat :=
  mul (ofInt 4) (square x)

def sourceRightDenominator (x : ExactRat) : ExactRat :=
  square x

def sourceLeft (x : ExactRat) (h : nonzero (sourceLeftDenominator x)) :
    ExactRat :=
  div (sourceNumerator x) (sourceLeftDenominator x) h

def claimedRight (x : ExactRat) (h : nonzero (sourceRightDenominator x)) :
    ExactRat :=
  div (claimedNumerator x) (sourceRightDenominator x) h

def witnessX : ExactRat := ofInt 1
def threeFourths : ExactRat := fraction 3 4 (by decide)
def three : ExactRat := ofInt 3

def witnessLeftDefined : nonzero (sourceLeftDenominator witnessX) := by
  decide

def witnessRightDefined : nonzero (sourceRightDenominator witnessX) := by
  decide

def witnessLeft : ExactRat := sourceLeft witnessX witnessLeftDefined
def witnessRight : ExactRat := claimedRight witnessX witnessRightDefined

/- This is the literal published final implication at the verified witness,
including x != 0, both division side conditions, and the valid numerator
identity. -/
abbrev PublishedFinalEqualityAtWitness : Prop :=
  nonzero witnessX ->
  nonzero (sourceLeftDenominator witnessX) ->
  nonzero (sourceRightDenominator witnessX) ->
  equivalent (sourceNumerator witnessX) (claimedNumerator witnessX) ->
  equivalent witnessLeft witnessRight

/-- Load-bearing refutation of the exact source equality at `x = 1`. It
certifies the domain premise, both defined divisions, the material numerator
identity, both exact values, the unequal conclusion, and the negated published
witness implication. -/
theorem reddit_learnmath_t6wb5n_refutation :
    nonzero witnessX /\
    nonzero (sourceLeftDenominator witnessX) /\
    nonzero (sourceRightDenominator witnessX) /\
    equivalent (sourceNumerator witnessX) (claimedNumerator witnessX) /\
    equivalent witnessLeft threeFourths /\
    equivalent witnessRight three /\
    Not (equivalent witnessLeft witnessRight) /\
    Not PublishedFinalEqualityAtWitness := by
  decide

def contextX : ExactRat := fraction 3 2 (by decide)
def epsilon : ExactRat := ofInt 1
def delta : ExactRat := fraction 4 5 (by decide)
def fourFifths : ExactRat := fraction 4 5 (by decide)
def sevenThirtySixths : ExactRat := fraction 7 36 (by decide)
def sevenNinths : ExactRat := fraction 7 9 (by decide)
def contextDistance : ExactRat := absRat (sub contextX (ofInt 2))

def contextLeftDefined : nonzero (sourceLeftDenominator contextX) := by
  decide

def contextRightDefined : nonzero (sourceRightDenominator contextX) := by
  decide

def contextLeft : ExactRat := sourceLeft contextX contextLeftDefined
def contextRight : ExactRat := claimedRight contextX contextRightDefined

/-- Context-strengthened rational real witness satisfying every stated
epsilon-delta neighborhood condition. It shows that the same denominator
error remains visible near `2`. -/
theorem context_strengthened_witness_inside_proof_neighborhood :
    less (ofInt 0) epsilon /\
    less (ofInt 0) delta /\
    equivalent delta (minRat (ofInt 1) (mul fourFifths epsilon)) /\
    less (ofInt 0) contextDistance /\
    less contextDistance delta /\
    nonzero contextX /\
    nonzero (sourceLeftDenominator contextX) /\
    nonzero (sourceRightDenominator contextX) /\
    equivalent (sourceNumerator contextX) (claimedNumerator contextX) /\
    equivalent contextLeft sevenThirtySixths /\
    equivalent contextRight sevenNinths /\
    Not (equivalent contextLeft contextRight) := by
  decide

def controlX : ExactRat := ofInt 2

def controlLeftDefined : nonzero (sourceLeftDenominator controlX) := by
  decide

def controlRightDefined : nonzero (sourceRightDenominator controlX) := by
  decide

def controlLeft : ExactRat := sourceLeft controlX controlLeftDefined
def controlRight : ExactRat := claimedRight controlX controlRightDefined

/-- Meaningful negative control: at `x = 2`, every domain and definedness
premise remains true, while the zero numerator makes both encoded source sides
zero and equal, masking precisely the missing denominator factor. -/
theorem negative_control_zero_numerator_masks_missing_factor :
    nonzero controlX /\
    nonzero (sourceLeftDenominator controlX) /\
    nonzero (sourceRightDenominator controlX) /\
    equivalent (sourceNumerator controlX) (ofInt 0) /\
    equivalent (claimedNumerator controlX) (ofInt 0) /\
    equivalent controlLeft (ofInt 0) /\
    equivalent controlRight (ofInt 0) /\
    equivalent controlLeft controlRight := by
  decide

end RedditLearnmathT6wb5n

\end{Vlean}
\subsection*{Proof-safety information}
The automated scan found no unfinished proofs, new axiom declarations, unsafe declarations, external code hooks, or uses of Lean's compiled decision procedure. This scan concerns proof construction only; it does not establish that the theorem matches the source claim.
\textit{Technical scan counts:} native\_decide=0, sorry=0, axiom\_decl=0, unsafe=0, extern=0. \textit{Technical tier:} KERNEL-ONLY.
\textbf{A YES means:} Lean represents the same objects and assumptions, uses the same proposed counterexample, and the final theorem statement asserts the same contradiction. Successful compilation by itself is not enough.
\clearpage
\section*{Human sign-off}
Complete this block from your own review of Questions 1--3. Do not use the automated appendix below as an answer key.\par
\textbf{Reviewer name:} \rule{9cm}{0.4pt}\par
\textbf{Date:} \rule{4cm}{0.4pt}\par
\medskip\textbf{1. Does the proof verification document accurately represent the source and challenged claim?}\par
\answerchoices\par
\decisionexplanation
\medskip\textbf{2. Does the proposed counterexample satisfy the assumptions and make that claim false?}\par
\answerchoices\par
\decisionexplanation
\medskip\textbf{3. Does the final Lean theorem statement represent and verify that same counterexample?}\par
\answerchoices\par
\decisionexplanation
\medskip\textbf{Overall decision:}\par
\checkbox\ VALIDATED (all three answers are YES)\par
\checkbox\ REJECTED (at least one answer is NO)\par
\checkbox\ NEEDS ANOTHER REVIEWER (at least one CANNOT DETERMINE and no NO)\par
\medskip\textbf{Additional comments.} Use this space for any extra source, counterexample, calculation, or Lean details that support your decisions.\par\rule{\linewidth}{0.4pt}\par\vspace{7pt}\rule{\linewidth}{0.4pt}
\clearpage
\section*{Appendix --- Prior automated review (background only)}
\textbf{These results are not an answer key.} They come from the earlier Opus/Fable verification campaign and are printed only so an auditor can see what was previously checked. The human decisions above must be made independently. Raw status codes are retained in small type for traceability.
\textbf{Overall automated disposition:} Cleared by automated screening; an independent search also found a counterexample.
\par\textit{Audit code:} \nolinkurl{PASS_CORROBORATED}
\subsection*{Independent search with no access to the stored example}
\textbf{Result:} An independent search found a counterexample without seeing the stored example.
\par\textit{Audit code:} \nolinkurl{WITNESS_FOUND}
\par\textbf{Cases checked:}
\begin{Vsrc}
Exhaustive over x = p/q in lowest terms with |p| <= 60, 1 <= q <= 60, p!= 0 (4406 distinct nonzero rationals), plus a targeted re-sweep of the 732 of those lying in the punctured neighbourhood 1 < x < 3, x!= 2 (the only regime the epsilon-delta proof uses). All 15 pairwise equalities among E1..E6 tested at each x.
\end{Vsrc}
\subsection*{Independent re-check of the proposed counterexample}
\textbf{Result:} A separate checker recomputed the proposed counterexample and confirmed it.
\par\textit{Audit code:} \nolinkurl{CONFIRMED}
\par\textbf{Checker's note:}
\begin{Vsrc}
Independent from-scratch recomputation reproduces the recorded witness exactly. At x = 1 the first four equalities of the quoted chain all hold (every term equals 3/4), and the fifth and final asserted equality fails: |(x-2)(x+2)|/(4x^2) = 3/4 whereas |x-2||x+2|/x^2 = 3, and 3/4!= 3. The defect is the silent loss of the factor 4 in the denominator, so the two sides differ by exactly a factor of 4 whenever the numerator is nonzero. (Shortened here; the complete note remains in the machine-check JSON.)
\end{Vsrc}
\subsection*{Independent review of the source-to-Lean correspondence}
\textbf{Result:} A separate reviewer judged that Lean proves the same failure using the same example.
\par\textit{Audit code:} \nolinkurl{MATCHES}
\par\textbf{Theorem summary from that reviewer:}
\begin{Vsrc}
In a zero-import exact-rational model (ExactRat: Int numerator, positive Nat denominator, equality by cross multiplication), the main theorem reddit_learnmath_t6wb5n_refutation is a `by decide` conjunction at witnessX = ofInt 1 asserting: x is nonzero; both divisors 4*x^2 and x^2 are nonzero (both divisions defined); the valid numerator identity |(x-2)(x+2)| = |x-2|*|x+2| holds at x=1; sourceLeft = |(x-2)(x+2)|/(4*x^2) is exactly 3/4; claimedRight = |x-2|*|x+2|/(x^2) is exactly 3; (Shortened here; the complete note remains in the machine-check JSON.)
\end{Vsrc}
\par\textbf{Difference noted by that reviewer:}
\begin{Vsrc}
None material. (1) Encoding fidelity: sourceLeftDenominator = mul (ofInt 4) (square x) retains the material factor 4 and sourceRightDenominator = square x, so the two encoded sides are literally the quoted final equality's LHS and RHS with no weakening or strengthening; (Shortened here; the complete note remains in the machine-check JSON.)
\end{Vsrc}
\subsection*{Follow-up review}
No follow-up review was needed or recorded.
\textit{This appendix is only background. You do not need any separate machine-output files to answer the three questions.}
\end{document}